\documentclass[12pt]{article}
\usepackage[T1]{fontenc}
\usepackage[utf8]{inputenc}
\usepackage{lmodern}
\usepackage{textcomp}
\usepackage{enumitem}
\usepackage{geometry}
\geometry{margin=1in}
\begin{document}
\begin{center}
Answer Key - Business Administration - Undergraduate Level Assessment 4 \
\small (Answers are ordered exactly as in the question paper.)
\end{center}

\section*{Multiple Choice}
\begin{enumerate}[label=\arabic*.]
\item \textbf{Answer:} D
\\ \textit{Options:} A) Private services, social, economic and environmental; B) Public services, social, cultural and environmental; C) Private services, cultural, economic and environmental; D) Public services, social, economic and environmental
\\ \textit{Note:} The UK Public Services (Social Value) Act 2013 mandates that public authorities, when procuring public services, must take into account the social, economic, and environmental benefits that can be achieved, thereby highlighting the importance of integrating broader societal impacts into procurement decisions.
\item \textbf{Answer:} D
\\ \textit{Options:} A) Authority; B) Regulations; C) Command structure; D) Change
\\ \textit{Note:} Change is not a characteristic of a bureaucratic organization because such organizations are typically defined by their rigid structures, established procedures, and resistance to flexibility, which contrasts with the dynamic nature of change.
\item \textbf{Answer:} C
\\ \textit{Options:} A) \$7,140; B) \$7,000; C) \$6,860; D) \$6,300
\\ \textit{Note:} The correct answer is $6,860 because the store took advantage of the 2% discount offered for paying within 10 days, reducing the total invoice amount of $7,000 by $140 (2% of $7,000), resulting in a payment of \$6,860.
\item \textbf{Answer:} D
\\ \textit{Options:} A) Executive; B) Functional; C) Line management; D) Supplier
\\ \textit{Note:} The term "supplier" refers to a business entity that provides goods or services, rather than a recognized type of authority, which typically includes categories such as legal, traditional, or charismatic authority in organizational theory.
\item \textbf{Answer:} B
\\ \textit{Options:} A) Defining the program; B) Planning the program; C) Taking action and implementing ideas; D) Evaluation of the program
\\ \textit{Note:} The correct answer is "Planning the program" because this stage involves establishing clear objectives that guide the direction and strategy of the overall plan.
\end{enumerate}

\section*{True / False}
\begin{enumerate}[label=\arabic*.]
\setcounter{enumi}{5}
\item \textbf{Answer:} False
\\ \textit{Options:} A) True; B) False
\\ \textit{Note:} Engagement encompasses the active interaction and communication between a business and its audience, which includes considering feedback and preferences to build meaningful relationships rather than merely advertising products.
\item \textbf{Answer:} False
\\ \textit{Options:} A) True; B) False
\\ \textit{Note:} The Ohio Studies into leadership primarily emphasize the behaviors and styles of leaders rather than solely focusing on starting and end positions.
\item \textbf{Answer:} True
\\ \textit{Options:} A) True; B) False
\\ \textit{Note:} The statement is true because, in digital marketing, the nature of disclosure is often inherently integrated into practices such as data collection and user consent, which can minimize the need for marketers to explicitly address disclosure as a separate issue.
\item \textbf{Answer:} True
\\ \textit{Options:} A) True; B) False
\\ \textit{Note:} Stakeholders often focus on a business's social, environmental, and ethical practices rather than solely on its financial performance when evaluating its broader impact on society.
\item \textbf{Answer:} False
\\ \textit{Options:} A) True; B) False
\\ \textit{Note:} McGregor's Theory X posits that workers are inherently lazy and require close supervision and direction, indicating that they do not thrive on significant responsibility and autonomy.
\end{enumerate}

\section*{Long Form Response}
\begin{enumerate}[label=\arabic*.]
\setcounter{enumi}{10}
\item \textbf{Answer:} Geert Hofstede's concept of power distance refers to the extent to which less powerful members of a society defer to more powerful members, highlighting how authority and inequality are perceived in different cultures. In high power distance cultures, hierarchical structures are prevalent, and there is an acceptance of unequal power distribution, which can lead to centralized decision-making in businesses. Conversely, low power distance cultures favor egalitarianism and encourage participative management styles. Understanding power distance is crucial for international business practices as it affects leadership styles, negotiation strategies, and employee interactions. Additionally, it informs cross-cultural communication by highlighting potential misunderstandings arising from differing expectations regarding authority and respect.
\\ \textit{Note:} The answer correctly identifies and contrasts high and low power distance cultures, illustrating how these cultural dimensions affect organizational hierarchy, decision-making processes, and communication styles in international business contexts.
\item \textbf{Answer:} Issues management plays a crucial role in public relations by enabling organizations to proactively identify, analyze, and respond to emerging issues that could affect their reputation and operations. This strategic approach involves monitoring the external environment and engaging with stakeholders to understand potential challenges and opportunities. By effectively managing issues, organizations can craft informed responses that mitigate negative impacts and capitalize on positive developments. Furthermore, a well-structured issues management process fosters organizational resilience and adaptability, ensuring that responses are timely and appropriate in the face of evolving circumstances. Ultimately, this proactive stance not only protects the organization's image but also enhances stakeholder trust and loyalty.
\\ \textit{Note:} correct because it emphasizes the proactive nature of issues management in public relations, highlighting its importance in identifying and analyzing emerging issues to shape informed organizational responses that protect reputation and leverage opportunities.
\end{enumerate}

\vfill
\noindent\textit{End of Answer Key}
\end{document}